% Richardson extrapolation of the DST eigenvalue sequence on the L-shaped
% domain: methodology, validation against the published ground state, and the
% limits of the method.
%
% A note accompanying results/eigenvalues_dof_requirement/, standalone: compile
% with pdflatex. The data quoted are those of
% results_paper/eigenvalues_convergence/L_shaped_convergence_lambda{1,1000}.csv
% and results/eigenvalues_dof_requirement/L_shaped_dof_requirement_n1000.csv.
%
% Generated by Claude Code (Claude Opus 5).

\documentclass[11pt,a4paper]{article}

\usepackage[T1]{fontenc}
\usepackage{amsmath}
\usepackage{amssymb}
\usepackage{booktabs}
\usepackage[margin=3cm]{geometry}

\newcommand{\dofs}{\mathrm{dofs}}

\title{Richardson extrapolation of the DST eigenvalue sequence\\
       on the L-shaped domain}
\author{}
\date{}

\begin{document}

\maketitle

\section{The problem}

The degree-of-freedom requirement computed in this directory measures each
discretisation against a reference set of eigenvalues. For the L-shaped domain
no such set is published: the only Dirichlet eigenvalue of the L-shape for which
a value of many digits exists is the ground state,
\begin{equation}
  \lambda_1 = 9.6397238440219,
  \label{eq:mps}
\end{equation}
computed by Betcke and Trefethen with the method of particular solutions. Away
from it the computation falls back on its own finest run, the DST spectrum at
$M = 163$, and the accuracy of that run is then a floor: a tolerance of the
order of its own error cannot be resolved, and near the floor the DST rows of
the answer table measure the agreement of two DST runs rather than the distance
of either from the exact eigenvalues.

The floor is set by the ground state and not, as one might expect, by the top of
the requested block. At $M = 141$ the DST value of $\lambda_1$ is
$9.64855601236$, which is $9.2 \times 10^{-4}$ from~\eqref{eq:mps}, whereas the
same run places $\lambda_{1000}$ within about $1.2 \times 10^{-5}$ of where the
sequence is heading. The re-entrant corner is what holds the ground state back,
and refining the grid chips at that error slowly: the step from $14840$ to
$25025$ degrees of freedom, which is a factor $1.69$ in the count and $1.30$ in
$1/h$, would improve it by a factor of about $1.4$.

This note describes a different way of using the same runs, which improves the
ground state by four orders of magnitude at no computational cost, and reports
how far it can be trusted. It documents a method; the computation in this
directory does not use it, and takes the $M = 163$ run as its reference.

\section{The grid and the error model}

The domain is $(-1,1)^2 \setminus [0,1]^2$, embedded in the bounding box
$[-1,1]^2$ and sampled on a uniform grid with
\begin{equation}
  h = \frac{2}{M+1},
  \qquad
  \dofs(M) = \frac{3M^2}{4} - \frac{M}{2},
\end{equation}
where $M$ is the number of interior points across one slice of the box. $M+1$ is
kept even throughout, so that the re-entrant corner at the origin falls on a grid
line.

Near a corner of interior angle $\omega$ the eigenfunctions of the Dirichlet
Laplacian carry a singular part
\begin{equation}
  u(r,\theta) \sim r^{\pi/\omega} \sin\!\left(\frac{\pi\theta}{\omega}\right),
\end{equation}
and for the re-entrant corner of the L-shape, $\omega = 3\pi/2$, the exponent is
$\pi/\omega = 2/3$. Such a $u$ lies in $H^{s}$ for $s < 1 + 2/3$ and no better.
For a Galerkin method the eigenvalue error is bounded by the square of the
energy-norm error of the eigenfunction, so an energy error $O(h^{2/3})$ gives
\begin{equation}
  \lambda^{(h)} - \lambda = O\!\left(h^{4/3}\right).
  \label{eq:rate}
\end{equation}
The DST discretisation is not a Galerkin method and~\eqref{eq:rate} is not a
theorem for it, but it is what its sequence exhibits, to four digits, as
Section~\ref{sec:validation} shows. That is the whole content of the model used
below: for a fixed eigenvalue index $k$,
\begin{equation}
  \lambda_k^{(h)} = \lambda_k + C_k\,h^{p_k} + o\!\left(h^{p_k}\right),
  \label{eq:model}
\end{equation}
with one dominant error term whose rate $p_k$ is not assumed known.

\section{Three-point extrapolation with an unknown rate}

Let three resolutions $h_1 > h_2 > h_3$ give values $\lambda^{(1)},
\lambda^{(2)}, \lambda^{(3)}$ of one eigenvalue. Dropping the remainder
in~\eqref{eq:model} leaves three equations
\begin{equation}
  \lambda^{(i)} = \lambda_\infty + C\,h_i^{\,p},
  \qquad i = 1,2,3,
  \label{eq:three}
\end{equation}
in the three unknowns $\lambda_\infty$, $C$ and $p$. Subtracting them in pairs
eliminates $\lambda_\infty$,
\begin{equation}
  \lambda^{(1)} - \lambda^{(2)} = C\left(h_1^{\,p} - h_2^{\,p}\right),
  \qquad
  \lambda^{(2)} - \lambda^{(3)} = C\left(h_2^{\,p} - h_3^{\,p}\right),
\end{equation}
and dividing eliminates $C$, leaving one scalar equation for the rate,
\begin{equation}
  F(p) \;=\;
  \frac{\lambda^{(1)} - \lambda^{(2)}}{h_1^{\,p} - h_2^{\,p}}
  \;-\;
  \frac{\lambda^{(2)} - \lambda^{(3)}}{h_2^{\,p} - h_3^{\,p}}
  \;=\; 0 .
  \label{eq:scalar}
\end{equation}
Equation~\eqref{eq:scalar} is solved by bisection on $p \in [0.3, 4]$, a bracket
wide enough to hold any rate these discretisations produce and narrow enough to
exclude the spurious root at $p = 0$, where both denominators vanish. Two hundred
bisections take the interval far below the precision of the data; the cost is
nil, so there is no reason to be clever. The rate in hand, the remaining two
unknowns follow by back-substitution,
\begin{equation}
  C = \frac{\lambda^{(1)} - \lambda^{(2)}}{h_1^{\,p} - h_2^{\,p}},
  \qquad
  \lambda_\infty = \lambda^{(1)} - C\,h_1^{\,p},
  \label{eq:back}
\end{equation}
and $\lambda_\infty$ is the extrapolated eigenvalue.

For~\eqref{eq:scalar} to have a root the three values must be strictly monotone,
$C$ being of one sign; a sequence that has already reached its rounding plateau,
or that is not yet in the asymptotic regime, will fail this and must be
rejected rather than extrapolated.

When the resolutions happen to be geometrically spaced, $h_2 = r h_1$ and $h_3 =
r h_2$ with $r < 1$, equation~\eqref{eq:scalar} has the closed-form solution
\begin{equation}
  p = \frac{1}{\log(1/r)}
      \log\!\frac{\lambda^{(1)} - \lambda^{(2)}}{\lambda^{(2)} - \lambda^{(3)}},
\end{equation}
which is the textbook Richardson estimate of an unknown order. The resolutions
of the sweep in this directory were chosen to give round degree-of-freedom
counts and are not geometrically spaced, so~\eqref{eq:scalar} is solved
numerically instead.

\section{Validation on the ground state}
\label{sec:validation}

The ground state is the one place where the extrapolation can be checked against
a value computed by an unrelated method. Table~\ref{tab:seq} is the DST sequence,
Table~\ref{tab:rich} what~\eqref{eq:scalar}--\eqref{eq:back} make of it.

\begin{table}[htbp]
  \centering
  \caption{The DST ground state on the L-shaped domain, and its distance from
           the value~\eqref{eq:mps} of the method of particular solutions.}
  \label{tab:seq}
  \begin{tabular}{rrrll}
    \toprule
    $M$ & $\dofs$ & $h$ & $\lambda_1^{(h)}$ & relative error \\
    \midrule
     75 &  4181 & 0.026315789 & 9.66004965197 & $2.11 \times 10^{-3}$ \\
     81 &  4880 & 0.024390244 & 9.65809106626 & $1.91 \times 10^{-3}$ \\
     93 &  6440 & 0.021276596 & 9.65503306574 & $1.59 \times 10^{-3}$ \\
    105 &  8216 & 0.018867925 & 9.65276691084 & $1.35 \times 10^{-3}$ \\
    115 &  9861 & 0.017241379 & 9.65128964236 & $1.20 \times 10^{-3}$ \\
    141 & 14840 & 0.014084507 & 9.64855601236 & $9.16 \times 10^{-4}$ \\
    \bottomrule
  \end{tabular}
\end{table}

\begin{table}[htbp]
  \centering
  \caption{Three-point extrapolation of the sequence of Table~\ref{tab:seq}, for
           several choices of the three resolutions. The rate $p$ is solved for,
           not imposed. The last column is the distance of $\lambda_\infty$
           from~\eqref{eq:mps}.}
  \label{tab:rich}
  \begin{tabular}{lrrll}
    \toprule
    resolutions & $p$ & $C$ & $\lambda_\infty$ & relative error \\
    \midrule
    $105, 115, 141$ & 1.3335 & 2.59775 & 9.63972466785 & $8.55 \times 10^{-8}$ \\
    $93, 115, 141$  & 1.3335 & 2.59791 & 9.63972480833 & $1.00 \times 10^{-7}$ \\
    $93, 105, 141$  & 1.3335 & 2.59803 & 9.63972492366 & $1.12 \times 10^{-7}$ \\
    $81, 105, 141$  & 1.3335 & 2.59828 & 9.63972515704 & $1.36 \times 10^{-7}$ \\
    $75, 93, 141$   & 1.3336 & 2.59866 & 9.63972556092 & $1.78 \times 10^{-7}$ \\
    $75, 81, 93$    & 1.3337 & 2.59965 & 9.63972724885 & $3.53 \times 10^{-7}$ \\
    \bottomrule
  \end{tabular}
\end{table}

Two things are worth reading off these tables. The first is the rate: every
triple returns $p = 1.3335$, against the $4/3 = 1.3333$ that the corner
regularity predicts in~\eqref{eq:rate}. The model behind the extrapolation is
therefore not a curve that happens to fit three points, but the mechanism that
actually limits the method, and the agreement is stable across triples spanning
the whole sweep.

The second is the gain. The finest run of the sequence is $9.2 \times 10^{-4}$
from the published value; the extrapolation of its three finest members is
$8.5 \times 10^{-8}$ from it, four orders of magnitude closer, at the cost of
some arithmetic on numbers already computed. Even the three coarsest runs of
Table~\ref{tab:seq}, which together carry fewer degrees of freedom than the
finest one alone, extrapolate to $3.5 \times 10^{-7}$.

\section{Away from the ground state}

Nothing outside the ground state can be checked against a published value, and
the behaviour there is not the same. Table~\ref{tab:deep} extrapolates
$\lambda_{1000}$ from the same runs.

\begin{table}[htbp]
  \centering
  \caption{Three-point extrapolation of $\lambda_{1000}$, whose DST values are
           $4363.82471980$, $4363.88743383$, $4363.92057849$ and
           $4363.96805292$ at $M = 93, 105, 115, 141$. There is no published
           value to compare with.}
  \label{tab:deep}
  \begin{tabular}{lrrl}
    \toprule
    resolutions & $p$ & $C$ & $\lambda_\infty$ \\
    \midrule
    $93, 105, 141$  & 3.2440 & $-51610$ & 4364.01901823 \\
    $93, 115, 141$  & 3.2104 & $-45572$ & 4364.01998755 \\
    $105, 115, 141$ & 3.1618 & $-37825$ & 4364.02107258 \\
    \bottomrule
  \end{tabular}
\end{table}

The rate is no longer $4/3$ and no longer stable: it drifts from $3.24$ to
$3.16$ as the triple moves, and the extrapolated values spread over
$2 \times 10^{-3}$ in absolute terms, that is $5 \times 10^{-7}$ relative. The
thousandth eigenvector oscillates on a scale much finer than the corner
singularity and is not limited by it, so the mechanism behind~\eqref{eq:rate}
does not apply; what governs the error instead is the resolution of the
oscillation, and the coarse runs of the sweep, which carry only a few thousand
degrees of freedom for a thousand eigenvalues, are barely in the asymptotic
regime at all. The large fitted $C$ is a symptom of the same thing: the model is
being stretched.

The extrapolation is still worth something there. The finest run sits
$1.2 \times 10^{-5}$ from where the three triples agree the sequence is heading,
and they agree among themselves twenty times more closely than that. But the
consistency of the ground state, where three digits of the rate are reproduced
and the answer is confirmed independently, is not available here, and a spread
across triples is the only diagnostic left.

\section{What the method does not give}

\begin{itemize}
  \item \textbf{No bound.} Extrapolation returns an estimate, not an interval.
        The spread over several triples is a diagnostic, not an error bar; it
        can be small for the same wrong reason in every triple.
  \item \textbf{One error term.} Equation~\eqref{eq:model} assumes a single
        dominant term. Where two terms are comparable, the fitted $p$ is an
        average of two rates and $\lambda_\infty$ inherits the discrepancy. This
        is what the drift of $p$ in Table~\ref{tab:deep} shows.
  \item \textbf{The asymptotic regime is required.} Runs must be fine enough for
        the remainder in~\eqref{eq:model} to be negligible against the leading
        term. For a high index $k$, that means many more than $k$ degrees of
        freedom, not merely more than $k$.
  \item \textbf{Rounding sets a limit.} The differences
        $\lambda^{(1)} - \lambda^{(2)}$ are formed by cancellation. For
        $\lambda \approx 4 \times 10^{3}$ in double precision, differences below
        roughly $10^{-9}$ are noise, and a sequence that has converged to
        rounding cannot be extrapolated at all --- it will usually fail the
        monotonicity test of~\eqref{eq:scalar} first.
  \item \textbf{Per-eigenvalue rates.} $p_k$ genuinely depends on $k$, as
        Tables~\ref{tab:rich} and~\ref{tab:deep} show. A reference set built this
        way must solve~\eqref{eq:scalar} for each $k$ separately and cannot
        impose one rate across the block.
\end{itemize}

\section{Reproduction}

The DST values are those of the run tables written by

\begin{itemize}
  \item \texttt{compute\_eigenvalues\_dof\_requirement.m},
        into this directory, and
  \item \texttt{make\_eigenvalues\_convergence\_index.m}, into\\
        \texttt{results\_paper/eigenvalues\_convergence/}.
\end{itemize}

The extrapolation is~\eqref{eq:scalar} with a bisection on $[0.3, 4]$ followed
by~\eqref{eq:back}. It needs nothing but the pairs $(h_i, \lambda^{(i)})$, and
no spectrum has to be recomputed for it.

\end{document}
